\DocumentMetadata{
  pdfversion=2.0,
  pdfstandard=ua-2,
  lang=en-US,
  tagging=on,
  tagging-setup={math/setup=mathml-SE,tabsorder=structure}
}
\documentclass[11pt,letterpaper]{article}
\usepackage[margin=0.68in,includefoot]{geometry}
\usepackage{newcomputermodern}
\usepackage{amsmath,amscd,mathtools}
\usepackage{tikz,adjustbox}
\providecommand{\square}{\mdlgwhtsquare}
\usepackage[hidelinks]{hyperref}
\hypersetup{
  pdftitle={Algebra PhD Exam, January 2018},
  pdfauthor={University of Florida Department of Mathematics},
  pdfsubject={Graduate mathematics examination},
  pdfdisplaydoctitle=true,
  bookmarksopen=true
}
\setcounter{secnumdepth}{0}
\setlength{\parindent}{0pt}
\setlength{\parskip}{0.28em}
\setlength{\emergencystretch}{3em}
\setlength{\leftmargini}{2.15em}
\setlength{\leftmarginii}{2.65em}
\setlength{\leftmarginiii}{3em}
\renewcommand{\labelenumi}{\textbf{\theenumi.}}
\pagestyle{empty}
\raggedbottom
\allowdisplaybreaks
\newenvironment{examproblems}[1]{%
  \begin{enumerate}%
  \setcounter{enumi}{#1}%
  \setlength{\topsep}{0.35em}%
  \setlength{\partopsep}{0pt}%
  \setlength{\itemsep}{0.48em}%
  \setlength{\parsep}{0.18em}%
}{\end{enumerate}}
\newenvironment{examsubparts}{%
  \begin{enumerate}%
  \setlength{\topsep}{0.18em}%
  \setlength{\partopsep}{0pt}%
  \setlength{\itemsep}{0.16em}%
  \setlength{\parsep}{0.1em}%
}{\end{enumerate}}
\newenvironment{examinstructions}{%
  \par\begingroup\itshape
}{%
  \par\endgroup\vspace{0.18em}
}
\makeatletter
\renewcommand\section{\@startsection{section}{1}{0pt}%
  {-1.25ex plus -.25ex minus -.1ex}%
  {0.55ex plus .1ex}%
  {\normalfont\large\bfseries}}
\renewcommand{\@maketitle}{%
  \newpage\null\vskip -1em%
  {\footnotesize\bfseries
    \href{https://www.ufl.edu/}{University of Florida}, \href{https://math.ufl.edu/}{Department of Mathematics}\par}%
  \vskip -0.5em%
  \tagstructbegin{tag=Title}%
    {\LARGE\bfseries\href{https://gma.math.ufl.edu/exams/algebra-phd/}{Algebra PhD Exam}, January 2018\par}%
  \tagstructend%
  \vskip 0.25em%
  \tagmcbegin{artifact}%
    \hrule height 1pt%
  \tagmcend%
  \vskip 1em%
}
\makeatother
\title{Algebra PhD Exam, January 2018}
\author{}
\date{}
\begin{document}
\maketitle
\thispagestyle{empty}
\begin{examinstructions}
Answer seven problems. If you turn in more, the first seven will be graded. Within reason, you may quote theorems as long as you state them clearly.
\end{examinstructions}
\begin{examinstructions}
Note. Below ring means associative ring with identity, and module means unital module.
\end{examinstructions}
\begin{examproblems}{0}
\item
Let \(q>1\) be a power of a prime~\(p\). Prove that there exists some field~\(F\) with \(|F|=q\).
\item
Determine the Galois group of \(x^4-5\) over~\(\mathbb{Q}\), over \(\mathbb{Q}(\sqrt{5})\), and over \(\mathbb{Q}(\sqrt{5}i)\). Justify your answers.
\item
Prove that if~\(G\) is a finite multiplicative subgroup of a field~\(F\), then~\(G\) is cyclic.
\item
Prove that if~\(R\) is an integral domain and a local ring, and~\(I\) an invertible ideal of~\(R\), then~\(I\) is principal.
\item
Let~\(R\) be a commutative ring with identity, and let~\(I\) and~\(J\) be ideals of~\(R\). Prove that the~\(R\)-modules \((R/I)\otimes_R(R/J)\) and \(R/(I+J)\) are isomorphic.
\item
State and prove Hilbert’s Basis Theorem.
\item
Prove the Lying-Over Theorem: Let~\(S\) be an integral extension of an integral domain~\(R\), and let~\(P\) be a prime ideal of~\(R\). Then, there exists a prime ideal~\(Q\) of~\(S\), such that \(Q\cap R=P\).
\item
Let~\(R\) be a ring.
\begin{examsubparts}
\item[\textnormal{(a)}]
Define what it means for an~\(R\)-module~\(M\) to be projective.
\item[\textnormal{(b)}]
Prove that any free~\(R\)-module is projective.
\item[\textnormal{(c)}]
Prove that there exists some ring~\(R\) and some projective~\(R\)-module~\(P\) such that~\(P\) is a projective module but not a free module.
\end{examsubparts}
\item
Let~\(R\) be a commutative ring with identity. Let \(\phi:R\to R\) be a ring homomorphism. Assume~\(R\) is Noetherian and that~\(\phi\) is surjective. Prove that~\(\phi\) is an automorphism.
\item
Let~\(V\) be a finite dimensional non-zero vector space over a field~\(F\). Let \(R=\operatorname{End}_F(V)\) be the ring of endomorphisms of~\(V\).
\begin{examsubparts}
\item[\textnormal{(a)}]
Prove that~\(R\) is a simple ring.
\item[\textnormal{(b)}]
Describe a minimal (non zero) left ideal of~\(R\).
\end{examsubparts}
\item
Let \(\mathcal{FG}\) be the category of all finite groups.
\begin{examsubparts}
\item[\textnormal{(a)}]
Recall the definition of free object in an arbitrary concrete category~\(\mathcal{C}\).
\item[\textnormal{(b)}]
Prove from your definition that there exists a finite set~\(S\) such that there is no free object in \(\mathcal{FG}\) freely generated by~\(S\).
\end{examsubparts}
\end{examproblems}
\end{document}
